% ============================================================================
% Fichier  : partie6/P06_C23_contre_test.tex
% Titre    : Le Contre-Test
% Auteur   : MINKA MI NGUIDJOÏ Thierry Emmanuel
% Version  : 1.0 -- novembre 2026
% ============================================================================
% Sources : pratiques de retest ; gestion des vulnérabilités. Contre-test de
% l'Opération LOGONE (novembre 2026) : cas fictif.
% ============================================================================

\chapter{Le Contre-Test}
\label{chap:retest}
\chaptermark{Le contre-test}

\epigraph{%
    « Une faille n'est corrigée que lorsqu'un second test le prouve. »%
}{Principe du cours}

\niveauChapitre{Intermédiaire}{Chapitre~\ref{chap:rapport}.}

\begin{boxobjectifs}
\begin{itemize}
    \item Définir le contre-test et sa place dans le cycle.
    \item Vérifier qu'une correction ferme réellement la faille.
    \item Distinguer correction de fond et palliatif.
    \item Documenter l'état de chaque constat après correction.
    \item Clore une mission proprement.
\end{itemize}
\end{boxobjectifs}

\bigskip

% ============================================================================
\section*{Question directrice : les corrections de LiZbiZ ont-elles vraiment fermé les failles~?}
% ============================================================================

Trois semaines après le rapport, M.~KOUNOU annonce que « tout est corrigé ».
NYONG SÉCURITÉ revient vérifier, constat par constat. Certaines failles sont
bien fermées, d'autres seulement masquées, une a été rouverte par une
nouvelle configuration. Le contre-test établit la vérité.

\medskip
\textbf{Plan du chapitre.}
\begin{enumerate}
    \item Qu'est-ce qu'un contre-test (section~\ref{sec:c23-definition}).
    \item Fond ou palliatif (section~\ref{sec:c23-fond}).
    \item L'état des constats (section~\ref{sec:c23-etats}).
    \item Clore la mission (section~\ref{sec:c23-cloture}).
\end{enumerate}

\separateur

% ============================================================================
\section{Qu'est-ce qu'un Contre-Test}
\label{sec:c23-definition}
% ============================================================================

\begin{boxdef}
\termecle{Contre-test}~: vérification, après correction, que chaque faiblesse
rapportée a bien été traitée, en rejouant le test qui l'avait révélée, dans
les mêmes conditions.
\end{boxdef}

% ============================================================================
\section{Fond ou Palliatif}
\label{sec:c23-fond}
% ============================================================================

\begin{boxalert}
\textbf{Violation typique.} Déclarer une faille corrigée parce qu'un symptôme
a disparu.

\smallskip
\textbf{Règle~:} on vérifie la cause, pas le symptôme. Masquer une bannière
de version ne corrige pas la version~; bloquer une adresse ne corrige pas le
défaut d'accès. On distingue clairement la correction de fond du palliatif
temporaire.
\end{boxalert}

% ============================================================================
\section{L'État des Constats}
\label{sec:c23-etats}
% ============================================================================

\begin{boxscenario}{Opération LOGONE. Le contre-test de novembre}
\textbf{Constatations.} Sur les 14~constats~: 9~fermés (dont T-01 et T-02,
critiques), 3~partiellement traités (T-07 corrigé sur les serveurs mais pas
les postes~; T-08 en cours~; T-11, sensibilisation engagée), 1~palliatif
(T-13, bannières masquées mais versions inchangées), 1~rouvert (T-10, un
nouveau service de test exposé).

\medskip
\textbf{Analyse} (confiance élevée). Les deux failles critiques sont bien
fermées. Le constat rouvert montre l'intérêt d'un contre-test~: sans lui, la
réexposition serait passée inaperçue.
\end{boxscenario}

\begin{tableaucours}{L{3.2cm}L{9.0cm}}{%
    \tableauHeader{État} &
    \tableauHeader{Signification}
}
Fermé & La faille ne se reproduit plus~; correction de fond vérifiée \\
Partiel & Une partie du périmètre reste exposée \\
Palliatif & Symptôme masqué, cause non traitée \\
Rouvert & La faille, fermée, est réapparue \\
Non traité & Aucune correction \\
\end{tableaucours}
\vspace{-0.8em}
\begin{center}{\small\itshape Tableau~23.1. États d'un constat au contre-test}\end{center}

% ============================================================================
\section{Clore la Mission}
\label{sec:c23-cloture}
% ============================================================================

\begin{boxprocedure}
\textbf{Clore une mission}
\begin{enumerate}
    \item Rejouer chaque test et noter l'état du constat.
    \item Rédiger un addendum au rapport avec les états.
    \item Distinguer fond et palliatif pour chaque correction.
    \item Convenir d'un suivi pour les constats partiels.
    \item Détruire les données de test et restituer les accès.
\end{enumerate}
\end{boxprocedure}

\begin{boxmetier}{Auditeur de sécurité et RSSI}
\textbf{Le contre-test ferme la boucle.} Sans contre-test, un rapport reste
une liste de promesses. Le RSSI inscrit le contre-test au contrat et en fait
la condition de clôture.
\end{boxmetier}

\separateur

\begin{boxresume}
\textbf{Ce qu'il faut retenir de ce chapitre~:}
\begin{itemize}
    \item \textbf{Le contre-test} rejoue le test après correction.
    \item \textbf{On vérifie la cause}, pas le symptôme.
    \item \textbf{Chaque constat} reçoit un état~: fermé, partiel, palliatif,
          rouvert, non traité.
    \item \textbf{Un constat rouvert} justifie à lui seul le contre-test.
    \item \textbf{La clôture} documente les états et détruit les données.
\end{itemize}
\end{boxresume}

\begin{boxinfo}
\textbf{Points de contrôle de l'élève.} Je sais~: définir le contre-test~;
distinguer fond et palliatif~; attribuer un état à un constat~; rédiger un
addendum~; clore une mission.
\end{boxinfo}

\bigskip

\begin{boxexo}[title={\ding{45}\quad Exercice~23.1. Définition \niveaubase}]
Qu'est-ce qu'un contre-test~?
\end{boxexo}

\begin{boxexo}[title={\ding{45}\quad Exercice~23.2. Fond ou palliatif \niveaubase}]
Masquer une bannière de version est-il une correction de fond ou un
palliatif~? Pourquoi~?
\end{boxexo}

\begin{boxexo}[title={\ding{45}\quad Exercice~23.3. États \niveaubase}]
Donnez les cinq états possibles d'un constat au contre-test.
\end{boxexo}

\begin{boxexo}[title={\ding{45}\quad Exercice~23.4. T-10 rouvert \niveaumoyen}]
Comment un constat fermé peut-il être rouvert~? Que recommander~?
\end{boxexo}

\begin{boxexo}[title={\ding{45}\quad Exercice~23.5. Partiel \niveaumoyen}]
T-07 est corrigé sur les serveurs mais pas les postes. Quel état, et quel
suivi~?
\end{boxexo}

\begin{boxexo}[title={\ding{45}\quad Exercice~23.6. Addendum \niveaumoyen}]
Rédigez l'addendum de contre-test pour T-01 et T-13.
\end{boxexo}

\begin{boxexo}[title={\ding{45}\quad Exercice~23.7. Condition de clôture \niveaumoyen}]
Pourquoi faire du contre-test une condition de clôture au contrat~?
\end{boxexo}

\begin{boxexo}[title={\ding{45}\quad Exercice~23.8. Bilan \niveaumoyen}]
Dressez le bilan chiffré du contre-test de LOGONE.
\end{boxexo}

\begin{boxexo}[title={\ding{45}\quad Exercice~23.9. Procédure \niveauexpert}]
Proposez une procédure de contre-test réutilisable pour les futures missions
de NYONG SÉCURITÉ.
\end{boxexo}

\begin{boxexo}[title={\ding{45}\quad Exercice~23.10. Indicateur \niveauexpert}]
Définissez un indicateur de qualité de correction, calculé à partir des
états de contre-test.
\end{boxexo}

\bigskip

\begin{boxinfo}
\textbf{Pour aller plus loin.}
\begin{itemize}
    \item PTES~\cite{ptes}.
\end{itemize}
\end{boxinfo}

\begin{boxinfo}
\textbf{Transition.} Lorsqu'une faille touche un produit utilisé par
d'autres, sa correction dépasse le client~: le Chapitre~\ref{chap:divulgation}
traite de la divulgation coordonnée.
\end{boxinfo}
